\documentclass[12pt, letterpaper]{article}

%-------Packages---------
\usepackage{amssymb,amsfonts}
\usepackage{mathptmx}
\usepackage{amsthm}
\usepackage{graphicx}
\usepackage[all,arc]{xy}
\usepackage{enumerate}
\usepackage{bbm}
\usepackage{dsfont}
\usepackage{mathrsfs}
\usepackage{tikz-cd}
\usepackage{setspace}
\usepackage{import}
\usepackage[utf8]{inputenc}
\usepackage{hyperref}
\usepackage{tikz}
\usepackage{float}
\usepackage[dvipsnames]{xcolor}
\usepackage[nocheck]{fancyhdr}
\usepackage[nointegrals]{wasysym}

\usepackage[left=1in,top=1in,right=1in,bottom=1in]{geometry}

\usepackage{mathtools}
\usepackage{setspace}
\usepackage{cleveref}
\usepackage{multicol}
\usepackage{mathpazo}
\usepackage{tcolorbox}
\tcbuselibrary{theorems,skins,breakable}
\usepackage{xparse}


\usepackage{xcolor, ifthen, tcolorbox}
\tcbuselibrary{theorems,skins,breakable}
% Theorem System
% The following environments are provided:
%   New Problem:        \begin{problem} ...
%   New Question Header:   \qh (then followed by subproblems)
%   Subproblem:     \begin{subproblem} ...
%   Problem Proof:  \begin{problemproof} ...
%   Proof:          \\begin{pf}{color} ...
%   Remark:         \rmk (sentence), \rmkb (block)

% Define custom colors
\definecolor{thmscol}{HTML}{F4565B} %For Problems
\definecolor{lemscol}{HTML}{FE844C} %For Lemmes
\definecolor{prpscol}{HTML}{00A7EF} %For proposition
\definecolor{corscol}{HTML}{23DAFF} %For corrolaries
\definecolor{rmkscol}{HTML}{6DEEC5} %For remarks
\definecolor{defscol}{HTML}{18C991} %For definitions
\definecolor{exmscol}{HTML}{7DB800} %For examples
\definecolor{clmscol}{HTML}{165c58} %For claims

%Change between dark(black)/light(white) mode
\newcommand{\colormode}{white}
\ifthenelse{\equal{\colormode}{white}}
      {\newcommand{\TextColor}{black}}
      {\newcommand{\TextColor}{white}}
\newcommand{\pbcolormain}{thmscol!80!\colormode}
\newcommand{\pbcolorsecond}{thmscol!38!\colormode}


% ============================
% Problem Counter
% ============================
\newcounter{subproblemcounter}
\newcounter{problemcounter}

\newcommand{\q}{
    \setcounter{subproblemcounter}{0}%
    \stepcounter{problemcounter} % Increment problem counter
    \color{\TextColor}Problem \arabic{problemcounter}: % Display the problem counter
}

\newcommand{\stepsub}{
    \stepcounter{subproblemcounter}%
    \color{\TextColor}(\alph{subproblemcounter})
}
% ============================
% Problem
% ============================
\newtcolorbox{pb}[1]{
  enhanced,
  frame hidden,
  titlerule=0mm,
  toptitle=1mm,
  bottomtitle=1mm,
  fonttitle=\bfseries\large,
  coltitle=black,
  colbacktitle=\pbcolormain,
  colback=\pbcolorsecond,
  title={\q #1},
  coltext=\TextColor
}

\newtcolorbox{spb}[1]{
  enhanced,
  frame hidden,
  titlerule=0mm,
  toptitle=1mm,
  bottomtitle=1mm,
  fonttitle=\bfseries\large,
  coltitle=black,
  colbacktitle=\pbcolormain,
  colback=\pbcolorsecond,
  title={\stepsub #1},
  coltext=\TextColor,
  attach boxed title to top left={xshift=3mm,yshift=-2mm},
  boxed title style={
    colback=\pbcolormain,
    colframe=\pbcolormain,
  }
}

\newtcolorbox{pbh}[1]{
    enhanced,
    frame hidden,
    titlerule=0mm,
    toptitle=1mm,
    bottomtitle=1mm,
    fonttitle=\bfseries\large,
    coltitle=black,
    colbacktitle=\pbcolormain,
    colback=\pbcolorsecond,
    title={\q #1},
    boxrule=0pt,
    interior hidden,
    attach boxed title to top left={xshift=3mm,yshift=-2mm},
    boxed title style={
        colback=\pbcolormain,
        colframe=\pbcolormain,
    }
}

\newenvironment{problem}[1]{%
  \newpage
  \begin{pb}{#1}
    }{
  \end{pb}
}

\newenvironment{subproblem}[1]{%
  \begin{spb}{#1}
    }{
  \end{spb}
}

\newenvironment{problemproof}{%
    \begin{pf}{\pbcolormain}
        }{
    \end{pf}
}

\newcommand{\qh}[1][]{
    \newpage
    \begin{pbh}{#1}\end{pbh} % Display the problem counter
}

% ============================
% Claim
% ============================

\newtcbtheorem[number within=problemcounter]{claim}{Claim}
{
    enhanced,
    frame hidden,
    titlerule=0mm,
    toptitle=1mm,
    bottomtitle=1mm,
    fonttitle=\bfseries\large,
    coltitle=black,
    colbacktitle=clmscol!50!\colormode,
    colback=clmscol!20!\colormode,
}{clm}

\NewDocumentCommand{\clm}{m+m}{
    \begin{claim*}{#1}{}
        #2
    \end{claim*}
}

\newenvironment{clmpf}{
	{\noindent{\it \textbf{Proof.}}
	\setlength{\parindent}{0pt}
	}
	\tcolorbox[blanker,breakable,left=5mm,parbox=false,
    before upper={\parindent15pt},
    after skip=10pt,
	borderline west={1mm}{0pt}{clmscol!50!\colormode}]
}{
    \textcolor{clmscol!50!\colormode}{\hbox{}\nobreak\hfill$\blacksquare$}
    \endtcolorbox
}

\NewDocumentCommand{\clmp}{m+m+m}{
    \begin{claim*}{#1}{}
        #2
    \end{claim*}

    \begin{clmpf}
        #3
    \end{clmpf}
}


% ============================
% Proof
% ============================

\newenvironment{pf}[1]{%
  \def\pfcolor{#1}% Store the color in a macro
  \noindent{\it \textbf{Proof.}}%
  \tcolorbox[blanker,breakable,left=5mm,parbox=false,%
    before upper={\parindent15pt},%
    after skip=10pt,%
    borderline west={1mm}{0pt}{\pfcolor},%
    coltext=\TextColor
  ]%
  \setlength{\parindent}{0pt}%
}{%
  \textcolor{\pfcolor}{\hbox{}\nobreak\hfill$\blacksquare$}%
  \endtcolorbox%
}

\newtcolorbox{solution}{
  blanker,
  breakable,
  left=5mm,
  parbox=false,%
  before upper={\parindent15pt},%
  after skip=10pt,%
  borderline west={1mm}{0pt}{\pbcolormain},%
  coltext=\TextColor
}


% ============================
% Remark
% ============================


\NewDocumentCommand{\rmk}{+m}{
    {\it \color{rmkscol!100!white}#1}
}

\newenvironment{remark}{
    \par
    \vspace{5pt}
    \begin{minipage}{\textwidth}
        {\par\noindent{\textbf{Remark.}}}
        \tcolorbox[blanker,breakable,left=5mm,
        before skip=10pt,after skip=10pt,
        borderline west={1mm}{0pt}{rmkscol!80!\colormode},
        coltext=\TextColor
        ]
}{
        \endtcolorbox
    \end{minipage}
    \vspace{5pt}
}

\NewDocumentCommand{\rmkb}{+m}{
    \begin{remark}
        #1
    \end{remark}
}

\usepackage[all]{xypic}

% Define a new command for problems
\renewcommand{\headrule}{\color{\TextColor}\hrulefill}
\newcommand{\C}{\mathbb{C}}
\newcommand{\R}{\mathbb{R}}
\newcommand{\Q}{\mathbb{Q}}
\newcommand{\Z}{\mathbb{Z}}
\newcommand{\N}{\mathbb{N}}
\renewcommand{\P}{\mathbb{P}}
\newcommand{\F}{\mathbb{F}}
\newcommand{\contr}{\Rightarrow \Leftarrow}
\newcommand{\s}{\(\,\)}
\renewcommand{\t}[1]{\text{#1}}
\newcommand{\ang}[1]{\left\langle #1 \right\rangle}
\newcommand{\coker}{\mathrm{coker}}

% Meta data
\setstretch{1.2}

\begin{document}
\color{\TextColor}
\pagecolor{\colormode}
\setlength{\parindent}{0pt}
\pagestyle{fancy}
\fancyhf{}
\fancyhead[LOE]{\color{\TextColor}\slshape{\textbf{Vincent Ferrara}}}
\fancyhead[ROE]{\color{\TextColor}\thepage}
\qh[Unique Factorization]
\begin{subproblem}{}
    In the ring $\Z[i]$, factor the element $17+9i$ into a product
of irreducibles.
\end{subproblem}
\begin{problemproof}
    $N(17+9i)=2 \cdot 5 \cdot 37$
    \begin{align*}
        \dfrac{(17+9i)(1-i)}{(1+i)(1-i)}=13-4i\\
        \dfrac{(13-4i)(2+i)}{(2+i)(2-i)}=6+i
    \end{align*}
    $17+9i=(1+i)(2-i)(6+i)$

    These are irreducible since their norms are primes.
\end{problemproof}

\begin{subproblem}{}
    In the ring $\Z[\omega]$, factor the elements $17+9\omega$ and
$17-9\omega$ into products of irreducibles.
\end{subproblem}
\begin{problemproof}
    $N(17-9\omega)=523$ which is prime

    Thus, $17-9\omega$ is irreducible.

    $N(17+9\omega)=7 \cdot 31$
    \begin{align*}
        \dfrac{(17+9\omega)(3+\omega)}{(3+\omega)(3+\omega^2)}=(6+5\omega)
    \end{align*}
    $17+9\omega=(3+\omega^2)(6+5\omega)$

    These are irreducible since their norms are primes.
\end{problemproof}

\qh[Ideal Factorization]
\addtocounter{subproblemcounter}{5}
\begin{subproblem}{}
    Show that $(7) = \mathfrak{p}_3 \mathfrak{p}_4$ for some maximal (hence prime
ideals) $\mathfrak{p}_3$ and $\mathfrak{p}_4$. Write down
what $\mathfrak{p}_3$ and $\mathfrak{p}_4$ are explicitly, i.e., identify generators for them.
\end{subproblem}
\begin{problemproof}
    First we know that $3^2=9 \equiv -5 \pmod{7}$.

    Thus, consider $(7,3-\sqrt{-5})(7,3+\sqrt{-5})=(49,21-7\sqrt{-5},21+7\sqrt{-5},14)=(7)$. 
    
    Since $49-3 \cdot 14=7$, and the rest are divisible by 7.

    Want to show that $(R /(7,3-\sqrt{-5})) \cong \Z/7\Z$

    Define: $\phi: R \to \Z/7\Z$ s.t. $\phi(a+b\sqrt{-5})=a+3b$.

    Homomorphism:

    $\phi(a+b\sqrt{-5}+c+d\sqrt{-5})=a+c+3(b+d)=\phi(a+b\sqrt{-5})+\phi(c+d\sqrt{-5})$

    $\phi((a+b\sqrt{-5})(c+d\sqrt{-5}))=\phi(ac+(ad+cb)\sqrt{-5}-5bd)=ac-5bd+3(ad+cb)=ac+9bd+3(ad+cb)=(a+3b)(c+3d)=\phi(a+b\sqrt{-5})\phi(c+d\sqrt{-5})$

    $\phi(1)=1$

    Surjective:

    Let $x \in \Z/7\Z$. $\phi(x)=x$

    Kernel:

    Let $x=a+b\sqrt{-5} \in \ker{\phi}$. $\phi(x)=0=a+3b$ 
    
    Thus, $a=-3b+7k$ for some $k \in \Z$.

    $x=-3b+7k+b\sqrt{-5}=7k+b(-3+\sqrt{-5}) \in (7,-3+\sqrt{-5})=(7,3-\sqrt{-5})$.

    Let $x \in (7,3-\sqrt{-5})$. Thus, $x=7a+b(3-\sqrt{-5})$

    $\phi(x)=\phi(7a)+\phi(3-\sqrt{-5})=0+3-3=0$.

    Thus, $R / (7,3-\sqrt{-5}) \cong \Z/7\Z$.

    Thus, $(7,3-\sqrt{-5})$ is a maximal ideal. Do something similar for $(7,3+\sqrt{-5})$
\end{problemproof}

\begin{subproblem}{}
    Show that $\mathfrak{p}_3$ and $\mathfrak{p}_4$ are not principal, but their
squares are principal.
\end{subproblem}
\begin{problemproof}
    Assume they are principal then there is an $a \in \Z[\sqrt{-5}]$ s.t. $(a)=(7,3-\sqrt{-5})$.

    Then $N(a) \mid 49$ and $N(a) \mid 14$. Thus, $N(a) \mid 7$. Thus, $N(a)=7$. Since $(7,3-\sqrt{-5})$ is proper.

    Thus, $7=c^2+5d^2$.

    Consider $7 \equiv c^2+5d^2 \pmod{5}$. 
    
    Thus, $7 \equiv c^2 \pmod{5}$. 
    
    Thus, $2 \equiv c^2 \pmod{5}$. There is no solution. 
    
    Thus, there is no integer solution to this equation. Thus, $a$ does not exist.

    Do a similar thing for the other ideal. Thus, we have show that $(7,3-\sqrt{-5}),(7,3+\sqrt{-5})$ are not principal.

    $(7,3-\sqrt{-5})(7,3-\sqrt{-5})=(49,7(3-\sqrt{-5}),4-6\sqrt{-5})$

    Consider $(2-3\sqrt{-5})$

    $49=(2+3\sqrt{-5})(2-3\sqrt{-5})$
    
    $4-6\sqrt{-5}=2(2-3\sqrt{-5})$

    $7(3-\sqrt{-5})=(3+\sqrt{-5})(2-3\sqrt{-5})$

    $2-3\sqrt{-5}=49-3(21-7\sqrt{-5})+4(4-6\sqrt{-5})$

    Thus, $(2-3\sqrt{-5})=(7,3-\sqrt{-5})^2$

    The same works for $(2+3\sqrt{-5})=(7,3+\sqrt{-5})^2$ since just conjugate all the equations above.
\end{problemproof}

\begin{subproblem}{}
    Write a computer program to determine which of the first hundred
odd primes remain prime in $R$. Do you notice a pattern? (Please find a partner to
work with on this problem.
Indicate in your answer the student that you worked with.)
\end{subproblem}
\begin{problemproof}
    Ethan Harr
    
    [11, 13, 17, 19, 31, 37, 53, 59, 71, 73, 79, 97, 113, 131, 137, 139, 151, 157, 173, 179, 191, 193, 197, 199, 211, 233, 239, 251, 257, 271, 277, 293, 311, 313, 317, 331, 337, 353, 359, 373, 379, 397, 419, 431, 433, 439, 457, 479, 491, 499]

    The pattern seems to be $11,13,17,19 \pmod{20}$
\end{problemproof}

\begin{problem}{Euclidean domains}
    Show that $\Z [\sqrt{-2}]$ is a Euclidean domain.
\end{problem}
\begin{problemproof}
    Let the norm function as our size function.

    Let $x=a+b\sqrt{-2},y=c+d\sqrt{-2} \in \Z[\sqrt{-2}]$ s.t. $y \neq 0$.

    Consider $x/y=e+f\sqrt{-2}$ s.t. $e,f \in \Q[\sqrt{-2}]$.

    Let $g,h \in \Z$ s.t. $g$ has a gap of less than or equal to one half away form $e$ and same for $h$ and $f$.

    Thus,

    \begin{align*}
        a+b\sqrt{-2}&=(c+d\sqrt{-2})(e+f\sqrt{-2})\\
        &=(c+d\sqrt{-2})(g+h\sqrt{-2}+(e-g)+(f-h)\sqrt{-2})\\
        &=(c+d\sqrt{-2})(g+h\sqrt{-2})+(c+d\sqrt{-2})((e-g)+(f-h)\sqrt{-2})
    \end{align*}

    $N((c+d\sqrt{-2})((e-g)+(f-h)\sqrt{-2}))=N((c+d\sqrt{-2}))N((e-g)+(f-h)\sqrt{-2})$

    $N((e-g)+(f-h)\sqrt{-2})=(e-g)^2+2(f-h)^2 \leq 1/4+1/2=3/4$

    Thus, $N(r)\leq3/4N(c+d\sqrt{-2})$

    Thus, $\Z[\sqrt{-2}]$ is a Euclidean domain.
\end{problemproof}
\end{document}